\chapter{Adaptive mesh refinement}
\label{section:amr}

\Peano\ realises a multi-step refinement and erasing\footnote{I try to speak of
refine and erase and not of refine and coarsen, as the latter is a term used in
a multigrid context and usually means that you still keep the fine data
somewhere. In contrast, \Peano's erase means that the grid is really thrown
away.} strategy.

\begin{enumerate}
  \item Each subpartition asks the action set (clone) for the refinement and
  coarsening decisions at the startup of the mesh traversal. For this, they
  query \texttt{getGridControlEvents()}.
  \item The mesh traversal uses this information to trigger erases or
  refinements.
  \item In the subsequent mesh traversal, the mesh is actually altered.
\end{enumerate}


\section{or-based refinement and vertex lifecycle}

\begin{remark}
 This is something I still have to write.
\end{remark}


\section{Multi-action set usage pattern}

In many \Peano\ applications, the refinement pattern is not known a priori.
Instead, the action sets build up the knowledge where to refine or coarsen:
They start with an empty \texttt{std::vector< peano4::grid::GridControlEvent >},
and push back events throughout the grid run-through.
As action sets are not persistent in-between mesh traversal sweeps, as codes
might switch between different action sets, and as \Peano\ clones one action set
per subdomain (again, the number of subdomains per rank can change over time),
we have to be careful with the data management. 


A popular pattern is to have a static \texttt{std::vector<
peano4::grid::GridControlEvent >}. In this case, an action set can build up a 
local set of control events, and commit them to the static set.
The static set is then returned by \texttt{getGridControlEvents()} which
consequently erases it to accept new events.


\begin{code}
  std::vector< peano4::grid::GridControlEvent >   _gridControlEvents;

  static std::vector< peano4::grid::GridControlEvent >   _activeGlobalGridControlEvents;
  static std::vector< peano4::grid::GridControlEvent >   _currentGlobalGridControlEvents;

  static tarch::multicore::BooleanSemaphore _semaphore;

public:
  static void activateNewSetOfGridControlEvents();
\end{code}

To make this work, you might want to work with two static sets: the currently
active one and a new one, as each subspacetree will query
\texttt{getGridControlEvents()} independently.
However, we can exploit the fact that each subtree/subdomain asks for the grid
control events before they actually trigger \texttt{beginIteration()}. 
So it is safe to hand out the static set and to clear it within the begin
iteration.


If you work with multiple action sets and an algorithm which switches between
the different ones, you will have to centralise the static control events.
One action set might feed into the event, while the next one delivers from this
centralised repository.
Alternatively, you can add another (static) boolean which you set to true
whenever \texttt{getGridControlEvents()} is invoked.
The central events are then cleared by \texttt{beginIteration()} if and only if
the boolean is set.
